\section{Pengantar Metode Quine-McCluskey}

Meskipun Peta Karnaugh merupakan alat yang sangat efektif untuk penyederhanaan fungsi Boolean, metode ini memiliki keterbatasan praktis yang signifikan terkait skalabilitas. K-Map bekerja dengan sangat baik untuk fungsi hingga 4 variabel, masih dapat digunakan (dengan kesulitan yang meningkat) untuk 5--6 variabel, tetapi menjadi tidak praktis untuk fungsi dengan lebih banyak variabel. Keterbatasan ini berasal dari fakta bahwa manusia sulit memvisualisasikan dan mengenali pola pada peta berdimensi tinggi \cite{rosen2019,saylor_cs202}.

\subsection{Keterbatasan K-Map untuk Variabel Banyak}

Tabel berikut merangkum pertumbuhan kompleksitas K-Map seiring bertambahnya jumlah variabel:

\begin{table}[htbp]
\centering
\renewcommand{\arraystretch}{1.3}
\begin{tabular}{|c|c|c|l|}
\hline
\textbf{Variabel} & \textbf{Jumlah Sel} & \textbf{Dimensi Grid} & \textbf{Keterangan} \\ \hline
2 & 4 & $2 \times 2$ & Sangat mudah \\ \hline
3 & 8 & $2 \times 4$ & Mudah \\ \hline
4 & 16 & $4 \times 4$ & Standar, efektif \\ \hline
5 & 32 & $4 \times 4 \times 2$ (dua peta) & Cukup sulit \\ \hline
6 & 64 & $4 \times 4 \times 4$ (empat peta) & Sangat sulit \\ \hline
7+ & 128+ & Tidak praktis & Diperlukan metode algoritmik \\ \hline
\end{tabular}
\caption{Skalabilitas Peta Karnaugh}
\label{tab:skalabilitas-kmap}
\end{table}

Untuk fungsi dengan 5 variabel, K-Map memerlukan dua buah peta 4 variabel yang harus dibandingkan secara tiga dimensi (sel-sel pada posisi yang sama di kedua peta juga dianggap bersebelahan). Untuk 6 variabel, diperlukan empat buah peta. Semakin banyak variabel, semakin sulit bagi manusia untuk mengidentifikasi semua kelompok potensial dan memastikan optimalitas solusi \cite{wiki_karnaugh}.

\subsection{Algoritma Quine-McCluskey}

\begin{definisi}[Algoritma Quine-McCluskey]
\logic{Algoritma Quine-McCluskey} (juga dikenal sebagai metode tabulasi) adalah prosedur algoritmik untuk menemukan semua prime implicant dari suatu fungsi Boolean dan kemudian memilih himpunan minimum prime implicant yang mencakup seluruh minterm. Metode ini dikembangkan oleh Willard V. Quine (1952) dan disempurnakan oleh Edward J. McCluskey (1956) \cite{wiki_quine_mccluskey}.
\end{definisi}

Algoritma Quine-McCluskey terdiri dari dua tahap utama. Tahap pertama adalah pencarian semua prime implicant melalui proses penggabungan sistematis. Minterm-minterm yang hanya berbeda pada satu bit digabungkan menjadi suku yang lebih sederhana (satu literal dieliminasi), dan proses ini diulang hingga tidak ada lagi penggabungan yang memungkinkan. Tahap kedua adalah pemilihan himpunan minimum prime implicant yang mencakup semua minterm, yang dilakukan menggunakan tabel cakupan (\textit{prime implicant chart}).

\begin{contoh}
\textbf{Ilustrasi Tahap 1: Penggabungan Minterm}

Diberikan fungsi $F(A,B,C) = \sum m(1, 2, 5, 6, 7)$.

\textbf{Langkah 1:} Kelompokkan minterm berdasarkan jumlah bit bernilai 1:
\begin{center}
\begin{tabular}{|l|l|l|}
\hline
\textbf{Grup} & \textbf{Minterm} & \textbf{Biner ($ABC$)} \\ \hline
Grup 1 (satu ``1'') & $m_1$, $m_2$ & 001, 010 \\ \hline
Grup 2 (dua ``1'') & $m_5$, $m_6$ & 101, 110 \\ \hline
Grup 3 (tiga ``1'') & $m_7$ & 111 \\ \hline
\end{tabular}
\end{center}

\textbf{Langkah 2:} Bandingkan minterm antar grup yang berurutan dan gabungkan yang berbeda 1 bit:
\begin{center}
\begin{tabular}{|l|l|l|}
\hline
\textbf{Pasangan} & \textbf{Biner} & \textbf{Hasil (--- = dieliminasi)} \\ \hline
$m_1, m_5$ & 001, 101 & $-01$ ($= B'C$) \\ \hline
$m_2, m_6$ & 010, 110 & $-10$ ($= BC'$) \\ \hline
$m_5, m_7$ & 101, 111 & $1-1$ ($= AC$) \\ \hline
$m_6, m_7$ & 110, 111 & $11-$ ($= AB$) \\ \hline
\end{tabular}
\end{center}

Minterm $m_1$ dan $m_2$ (dari Grup 1) tidak dapat digabungkan karena berbeda 2 bit. Prime implicant yang diperoleh: $B'C$, $BC'$, $AC$, $AB$.
\end{contoh}

\subsection{Perbandingan K-Map dan Quine-McCluskey}

Tabel~\ref{tab:kmap-vs-qm} merangkum perbandingan kedua metode penyederhanaan.

\begin{table}[htbp]
\centering
\renewcommand{\arraystretch}{1.3}
\begin{tabular}{|>{\raggedright\arraybackslash}p{4.5cm}|>{\raggedright\arraybackslash}p{4.5cm}|>{\raggedright\arraybackslash}p{4.5cm}|}
\hline
\textbf{Aspek} & \textbf{Peta Karnaugh} & \textbf{Quine-McCluskey} \\ \hline
Pendekatan & Visual, pengenalan pola & Algoritmik, tabulasi \\ \hline
Jumlah variabel efektif & 2--4 (maks 6) & Tidak terbatas (secara teori) \\ \hline
Jaminan optimalitas & Tidak (bergantung pengguna) & Ya (selalu menemukan solusi minimal) \\ \hline
Otomatisasi & Sulit diotomatisasi & Mudah diimplementasikan di komputer \\ \hline
Kecepatan untuk kasus kecil & Sangat cepat & Lebih lambat (overhead tabulasi) \\ \hline
Kompleksitas komputasi & --- & Eksponensial pada kasus terburuk \\ \hline
\end{tabular}
\caption{Perbandingan Peta Karnaugh dan Algoritma Quine-McCluskey}
\label{tab:kmap-vs-qm}
\end{table}

Dalam praktik perancangan digital modern, alat-alat EDA (\textit{Electronic Design Automation}) menggunakan varian algoritmik dari Quine-McCluskey beserta heuristik lainnya (seperti algoritma Espresso) untuk mengoptimasi rangkaian secara otomatis. Namun, pemahaman terhadap Peta Karnaugh tetap penting sebagai fondasi konseptual, karena memberikan intuisi visual tentang mengapa dan bagaimana penyederhanaan Boolean bekerja \cite{wiki_quine_mccluskey}.

\subsection{Jembatan ke Bab Selanjutnya}

Pada Bab~11, hasil penyederhanaan yang diperoleh melalui K-Map akan diimplementasikan dalam bentuk rangkaian gerbang logika fisik. Mahasiswa akan mempelajari bagaimana ekspresi Boolean yang telah disederhanakan diterjemahkan menjadi diagram rangkaian menggunakan simbol-simbol gerbang logika standar, serta bagaimana rangkaian tersebut dapat dioptimasi lebih lanjut menggunakan gerbang universal NAND atau NOR.
